\documentclass[12pt]{article}
\usepackage[utf8]{inputenc}
\usepackage[T1]{fontenc}
\usepackage{hyperref}
\title{Piscataway Language Resources for K-12 Curriculum Development}
\author{Compiled by Claude Sonnet 5.5 (claude-sonnet-5-5) for Andr\'as Kornai}
\date{}
\begin{document}
\maketitle

\section{Introduction}
Piscataway (ISO 639-3 \texttt{psy}; Glottocode pisc1239) is an extinct Eastern Algonquian language of the Algic family, once spoken in the Maryland Chesapeake region. It is very thinly documented, and some sources treat it as close to or a dialect of Nanticoke. No living speakers are known. Searches found no verified school curriculum, primer or community language programme for it; this report therefore lists only background and archival leads, and a K-12 heritage curriculum would have to be built largely from related Eastern Algonquian material and community consultation.

\section{Primary Educational Resources}
No primers, textbooks or classroom materials for Piscataway were found in search results.

\section{Linguistic and Reference Materials}
\begin{itemize}
\item Wikipedia, ``Piscataway language'': overview of classification, status and attestation. \url{https://en.wikipedia.org/wiki/Piscataway_language}
\item Endangered Languages Project entry for Piscataway, which the search summary linked to a descriptive analysis of the surviving catechism text. \url{https://www.endangeredlanguages.com/node/108781}
\item ASJP database wordlist for Piscataway (a short comparative wordlist; the search summary indicated it draws on Mackie's 2006 Oxford MA thesis, \emph{Fragments of Piscataway}, which I did not locate online). \url{https://asjp.clld.org/languages/PISCATAWAY}
\end{itemize}

\section{Religious and Community Resources}
\begin{itemize}
\item Wikipedia, ``Piscataway people'': background on the Piscataway Conoy and related present-day communities, useful as a starting point for contacting tribal organizations. No verified language programme from them was found. \url{https://en.wikipedia.org/wiki/Piscataway_people}
\end{itemize}
The main surviving text is reported to be a short Catholic catechism attributed to the Jesuit Andrew White, held at Georgetown University; I did not see a direct link to it.

\section{Digital and Audio Resources}
No audio, video or software resources were found.

\section{Curriculum Development Context}
With no speakers and only fragmentary records, any curriculum would be a reclamation project: word lists, the catechism text and comparison with documented relatives such as Nanticoke and Delaware. It should be developed with Piscataway tribal communities, who decide what is shared and taught. Content would center on place names, culture and history rather than full language instruction.

\end{document}